\documentclass[11pt,a4paper]{article}
\usepackage[a4paper,margin=25mm]{geometry}
\usepackage[T1]{fontenc}
\usepackage[utf8]{inputenc}
\usepackage[english]{babel}
\usepackage{lmodern,microtype}
\usepackage{amsmath,amssymb,amsthm}
\usepackage{booktabs,tabularx,array}
\newcolumntype{Y}{>{\raggedright\arraybackslash}X}
\newcolumntype{P}[1]{>{\raggedright\arraybackslash}p{#1}}
\usepackage{xcolor}
\usepackage[hidelinks]{hyperref}
\usepackage{enumitem}
\usepackage{parskip}
\usepackage{fancyhdr}
\usepackage{needspace}
\setlist[itemize]{nosep,topsep=3pt,leftmargin=1.6em}
\setlist[enumerate]{nosep,topsep=3pt,leftmargin=1.6em}
\renewcommand{\arraystretch}{1.19}
\definecolor{ink}{HTML}{263445}
\definecolor{soft}{HTML}{66758B}
\definecolor{accent}{HTML}{254768}
\hypersetup{pdftitle={Ungoverned Transduction: When AI-Enabled Effects Outrun Authority},pdfauthor={Henric T. Boehm},pdfsubject={Concept Note v0.3.1: release-hardening with unchanged synthetic model}}
\pagestyle{fancy}
\fancyhf{}
\fancyhead[L]{\scriptsize\color{soft} UNGOVERNED TRANSDUCTION \;|\; CONCEPT NOTE v0.3.1}
\fancyhead[R]{\scriptsize\color{soft} 27 OCT 2026}
\fancyfoot[C]{\small\thepage}
\renewcommand{\headrulewidth}{0.25pt}
\newcommand{\heading}[1]{\vspace{2pt}\textbf{\color{accent}#1}}
\newcommand{\NE}{\textsc{Non-Evaluable}}
\newcommand{\NAdm}{\textsc{Non-Admissible}}
\newcommand{\AD}{\textsc{Admissible}}
\begin{document}
\thispagestyle{plain}
\begin{center}
{\LARGE\bfseries\color{ink} Ungoverned Transduction}\\[5pt]
{\large When AI-Enabled Effects Outrun Authority}\\[10pt]
{\normalsize A Concept Note on Non-Transduction, Status Promotion, and the Output-to-Effect Boundary}\\[10pt]
{\normalsize Henric T. B\"ohm}\\[4pt]
{\small \textbf{Independent Concept Note v0.3.1} \;---\; 27 October 2026}
\end{center}
\vspace{3pt}
\noindent\textbf{Status and boundary.} Independently legible application-level concept note with a bounded, executable \emph{synthetic} countermodel. Its logical and computational findings are conditional on stipulated semantics. The note is not a theorem of Holomorphic Spiral Resonance Theory (HSRT), nor of the project-internal HSRT-Core specification; it is not a validated extension to any pre-action framework, a security standard, evidence of deployment failure, or proof of operational superiority. The review is AI-assisted and author-directed, not independently endorsed; release and integrity declarations require human sign-off.

\noindent\textbf{Abstract.} AI capability may reach no observable effect, or may cause an effect without a warranted transition into execution. This note asks whether distinguishing these outcomes adds diagnostic clarity and, separately, whether a reading informed by the earlier Pre-Action Differentiation Cascade (PADC) adds value beyond competent enterprise authorization. We define an effect-window classification, stage-indexed admissibility, and a three-valued guard assessment; then construct a procurement-status countermodel and execute eight deterministic synthetic test cases against an end-to-end authorization baseline. Within that stipulated model, local permissions coexist with a failed global mandate, an admissible effect can have negative value, and missing evidence remains non-evaluable. The same authority-promotion defect is detected by a sufficiently specified baseline, so the test establishes a possible failure pattern but \emph{not} PADC-specific superiority, novelty or prevalence. A frozen claim inventory, explicit defeat conditions and a prospective comparator protocol specify what further evidence would be required.

\noindent\textbf{Terms and version scope.} \emph{Output-to-Action Admissibility} (OTAA) names the output-to-action admission question. \emph{Pre-Action Differentiation Cascade} (PADC) v0.4.0 is the \emph{earlier source actually consulted} for intake and guard vocabulary \cite{padc2026}. This note does \emph{not} claim alignment with the later \emph{Pre-Action Differentiation and Control} baseline. \emph{Distinction, Execution, and Adjudication} (DEA) supplies only \emph{proposed} role questions; the authoritative DEA publication master has not been compared. These acronyms do not denote independently validated methods here.

\noindent\textbf{Keywords.} Ungoverned transduction; non-transduction; AI governance; OTAA; PADC; status-to-authority promotion; delegated authority; Shadow AI; workflow authorization; non-evaluability.

\section{Problem: two different boundaries after capability}
The discussion essay \emph{The Silent AI Bubble} argues that scalable compute accumulates AI capability without automatically generating enterprise value \cite{boehm2026essay}. That account focuses on the gap between available capability and responsible workflow integration. A further question emerges when AI systems already act: what if capability \emph{does} reach an effect, but the transition is not adequately governed?

\textbf{Research question.} Under a declared observation window, can failure to produce an effect be distinguished from the production of an effect without a warranted cross-boundary mandate? Does a PADC-oriented reading discriminate such failures any better than a strong ordinary authorization review?

\textbf{Analytical method.} This is a conceptual decomposition with stipulated, stage-indexed predicates and one deliberately minimal countermodel. We (i) define objects, transitions and decision rules; (ii) construct an existence witness and positive control; (iii) run a deterministic eight-case Python model; and (iv) compare its verdicts against an equally informed end-to-end policy-enforcement baseline. The model \emph{demonstrates consistency within its encoding}, not real-world frequency or causal performance.

\textbf{Contribution.} A typed vocabulary for locating authority-promotion errors, an executable witness and failure-oriented evaluation protocol. \textbf{Non-contribution.} No new access-control algorithm, proof of framework independence, measured detection gain, or economic causal account.

The proposed thesis is deliberately narrower than a theory of the AI investment cycle:
\begin{equation}
\boxed{\begin{gathered}
\text{AI capability can fail before effect,}\\[-2pt]
\text{or reach effect while the governance transition fails.}
\end{gathered}}
\label{eq:thesis}
\end{equation}
The distinction is not a causal explanation of EBIT, productivity, or market valuations. It separates two possible failure locations that may be empirically investigated under declared conditions.

\section{Terminological repair: effect, admissibility, and value}
The original essay uses \emph{transduction} for the \emph{guarded} conversion of capability into action. Calling an unguarded effect \emph{ungoverned transduction} would therefore be contradictory without an explicit change of vocabulary. This note adopts two levels:
\begin{itemize}
\item \textbf{Descriptive transduction:} an AI-mediated chain produces an identified change in an external or organizational state; whether this change was permitted is undecided.
\item \textbf{Admissible transduction:} the identified output-to-effect transition passes a declared, stakes-appropriate object, evidence, role, mandate, risk, and trace contract.
\item \textbf{Value realization:} the realized effect satisfies a separately declared benefit, cost, quality, and risk evaluation. Neither execution nor admissibility establishes positive net value.
\end{itemize}

Let $C\in\{0,1\}$ denote demonstrated capability \emph{invoked on an eligible object}; $W$ an adequate, predeclared observation window; $E_W\in\{1,0,\bot\}$ an effect observation; $G_e\in\{1,0,\bot\}$ the end-to-end execution guard verdict; and $V_W$ a separately specified, possibly multidimensional, outcome evaluation. In either three-valued domain, $1$ means established/observed, $0$ means established not to hold, and $\bot$ means \emph{not evaluable from the declared evidence}. The verdict labels for $G_e$ are $1=\mathrm{Adm}$ (admissible), $0=\mathrm{NAdm}$ (non-admissible), and $\bot=\mathrm{NE}$ (non-evaluable); \emph{NAdm} is distinct from the MPRS abbreviation N/A (not applicable). These labels are outputs of the decision rule, not values used in the conjunction.

The classification presupposes a declared invocation, effect criterion and sufficient window. For evaluable cases:
\begin{equation}
\begin{aligned}
(C=1)\land(E_W=0)&\quad\rightsquigarrow\quad \text{non-transduction in }W,\\
(E_W=1)\land(G_e=0)&\quad\rightsquigarrow\quad \text{ungoverned transduction},\\
(E_W=1)\land(G_e=1)&\quad\rightsquigarrow\quad \text{admissible transduction}.
\end{aligned}
\label{eq:classes}
\end{equation}
When $C=0$, the invocation criterion is unmet and the case is \emph{out of scope}. When $E_W=\bot$, the effect classification is \NE{}; when $E_W=1$ but $G_e=\bot$, the guard classification is \NE{}. The label \emph{non-transduction} applies only to an adequate window with an observed non-effect, and does not assert permanent impossibility.
Here ``ungoverned'' means \emph{a demonstrably non-admissible output-to-effect transition under a frozen, justified governance contract}, not simply the use of a non-inventoried tool. An observed effect with unknown mandate or incomplete trace is \NE\ for that guard claim, unless a known violation independently settles it. A separately demonstrated documentation breach may be found, but the breach must not be inferred merely from the reviewer's lack of access. \emph{Non-transduction} is descriptive: it also includes a correctly blocked unsafe attempt and does not, by itself, diagnose an enterprise value failure.

A fourth possibility matters: $(E_W=1)\land(G_e=1)$ while $V_W\leq 0$. Properly authorized actions can be ineffective or costly. Thus this note does not equate \emph{admissible} with \emph{valuable}.

\section{Typed transition predicates and composition}
For a proposed effect $e$, distinguish the status-intake decision $A_s$, execution mandate $M_e$, status-to-effect semantic compatibility $S_{s\to e}$, evidence sufficiency $Q_e$, and local API permissions $L$. Every predicate is evaluated relative to a stated object, principal, action, policy version, time and evidence surface. In the \emph{stipulated model}, the complete execution guard is
\begin{equation}
G_e := L\land A_s\land S_{s\to e}\land M_e\land Q_e.
\label{eq:guard}
\end{equation}
For three-valued observations in $\{1,0,\bot\}$, a known false operand makes $G_e=0$; otherwise a missing operand makes $G_e=\bot$; all true operands yield $G_e=1$. Equivalently, the guard decision is $1\mapsto\mathrm{Adm}$, $0\mapsto\mathrm{NAdm}$, and $\bot\mapsto\mathrm{NE}$. This is a \emph{declared rule}, not a universal account of valid governance.

The correct, \emph{typed} non-entailments are
\begin{equation}
\boxed{\mathrm{Visible}(s)\not\Rightarrow A_s,\qquad A_s\not\Rightarrow M_e,
\qquad M_e\not\Rightarrow V_W>0.}
\label{eq:nonequiv}
\end{equation}
In contrast, by the stipulated complete-guard definition, $G_e=1\Rightarrow M_e=1$. The untyped shorthand ``admissibility does not imply authority'' is therefore \emph{withdrawn}: it was ambiguous and could contradict a fully specified execution-admissibility contract. A permitted read, status update, or API invocation need not imply a permitted downstream procurement effect. The implication changes only when the actions and guard scopes coincide.

\textbf{Conditional composition proposition (application level).} Suppose the local contracts prove only a status-write authorization and a service-account API authorization: $L=g_1\land g_2$. If neither contract carries the \emph{promotion} relation $S_{s\to e}$ or the effect mandate $M_e$, then a model with
\begin{equation}
\boxed{L=1,\qquad S_{s\to e}=0,\qquad M_e=0,\qquad G_e=0}
\label{eq:composition}
\end{equation}
is consistent. This is an \emph{existence witness under stipulated assumptions}, not proof that all local guards or all real deployments overlook global authorization. The positive control sets the status semantics and effect mandate to true. A correct end-to-end enforcement point checks $G_e$ and blocks the negative case. Thus local-to-global failure is \emph{conditional on incomplete composition coverage}, not an intrinsic limitation of access control.

\section{Constructed counterexample: status silently becomes authority}
Consider a \emph{synthetic} enterprise procurement workflow. This is a deliberately weakened local-control configuration, \emph{not} a fully implemented governance reference architecture. An AI agent is registered, authenticated through a non-human identity, allowed to update a vendor-readiness field, and comprehensively telemetried. It writes \texttt{READY\_FOR\_REVIEW}. A downstream automation improperly treats this status as an order approval and issues a purchase order.

\begin{center}
\small
\begin{tabularx}{\linewidth}{@{}P{0.22\linewidth}Y P{0.25\linewidth}@{}}
\toprule
\textbf{Step} & \textbf{Declared operation} & \textbf{Open question}\\
\midrule
$T_1$: output $\to$ status & Agent assesses vendor information and writes \texttt{READY\_FOR\_REVIEW}. & Was this status warranted?\\
$T_2$: status $\to$ authority-bearing trigger & Workflow maps \texttt{READY\_FOR\_REVIEW} to an approval precondition. & Who authorized the semantic promotion?\\
$T_3$: trigger $\to$ effect & Order service sends a purchase order with an account permitted to call the API. & Was \emph{this order} mandated?\\
\bottomrule
\end{tabularx}
\end{center}

The critical defect is not necessarily missing authentication, inventory, logging, or explainability; those controls may be present. It is a \emph{semantic and normative mismatch} between the scope of the upstream status and the authority claimed downstream. Nor does the example show that the eight-layer Shadow-AI approach \emph{must} permit this behavior: correctly implemented access policy, meaningful human approval, or end-to-end workflow authorization could prevent it. The example probes \emph{partial local implementations}, rather than refuting the full framework.

Two candidate, action-scoped evaluation obligations are:
\begin{align}
\textit{Status must not outrun its warrant},& \label{eq:invariant1}\\
\textit{Effect authority must not outrun its mandate}.& \label{eq:invariant2}
\end{align}
These are proposed review questions, not independent mathematical laws. Under the frozen procurement contract, the status $\texttt{READY\_FOR\_REVIEW}$ does not entail an approved order. A successful diagnosis identifies a specific violated conversion or mandate rule; participation of AI alone cannot yield a negative verdict.

\section{External problem domain and disciplined positioning}
Shunmugasundaram \emph{et al.} present Shadow AI as a governance problem associated with decentralized deployment, hidden integrations, non-human identities, autonomous execution, explainability and accountability \cite{shun2026}. Their publisher abstract describes governance elements including discovery, telemetry, Zero Trust, explainability and continuous auditability; the eight-part outline used in the motivating comparison also includes inventory, non-human identity and human oversight. The publisher abstract and bibliographic record were checked; the complete article's detailed governance design was \emph{not independently reconstructed here}. The work already addresses authorization and oversight: it should not be described as visibility-only.

Its use here is that of an \emph{external mechanism candidate and adversarial application domain}. The proposed addition is a transition-indexed question: even when the AI actor, credentials and actions are visible, exactly \emph{which proposition, status or delegated permission} was accepted as sufficient warrant for the ensuing effect? No empirical claim is made that the paper documents the synthetic case, demonstrates a transduction theory, or causally explains a financial outcome.

The NIST AI Risk Management Framework treats governance and AI risks broadly \cite{nist2023}. NIST SP 800-207 describes resource-centred authentication and authorization in Zero Trust architectures \cite{rose2020}; NIST SP 800-162 describes attribute-based decisions over subject, object, operation and environment attributes \cite{hu2019}. A policy enforcement point plus an explicit workflow/approval mandate can already implement the present $G_e$ rule. The relevant test for OTAA/PADC is therefore \emph{incremental diagnostic value beyond this credible baseline}, not novelty of authorization itself.

\section{Provisional PADC/DEA and interface crosswalk}
The following mapping is a candidate test interface, not a statement that the source documents prove an isomorphism:
\begin{center}
\small
\begin{tabularx}{\linewidth}{@{}P{0.26\linewidth}P{0.32\linewidth}Y@{}}
\toprule
\textbf{Observed event} & \textbf{PADC/OTAA burden} & \textbf{DEA / authority question}\\
\midrule
AI recommendation accepted by employee & Output-to-reliance; sufficiency of object, source, and evidence & Who may adjudicate its relevance?\\
Agent writes status to record & Output-to-status; status semantics and trace & Who may certify or promote this status?\\
Status triggers purchase/order/API & Status-to-effect; mandate, delegated scope, risk & Who may authorize and execute this effect?\\
Delegated agent invokes another agent & Chain-of-delegation; scope-preserving composition & Where is authority attenuated or expanded?\\
\bottomrule
\end{tabularx}
\end{center}

The earlier \emph{Pre-Action Differentiation Cascade} v0.4.0 explicitly maintains a separation between exposure, belief and action, uses Source Intake Mode for external claims, and requires trace and missingness declarations \cite{padc2026}. The HSRT-Core Interface Specification v0.2.2 requires an object, evidence surface, trace rule and window before evaluating a claim; gaps yield \NE{} rather than a negative verdict \cite{interface2026}. The present note therefore remains an \emph{Application Note / candidate Interface Artifact}; it changes no Core state, dynamics, axioms or theorems. The current \emph{Pre-Action Differentiation and Control} baseline and the authoritative DEA publication master both require a separate concordance check before any formal adoption of this crosswalk.

\section{Finite witness, positive control and observed model results}
We froze a small Python 3 standard-library script, \path{UT_v0_2_finite_model.py}, with eight named cases. Each case declares invocation, effect observation and window adequacy; local permissions; status admission; semantic compatibility; effect mandate; evidence sufficiency; and whether a complete policy enforcement point is operative. The classifier first distinguishes out-of-scope or unevaluable effects from observed non-effects; only observed effects receive an admissibility verdict. \emph{Execution of the script} is the reported method result; no organizational data are used.

\begin{center}
\small
\begin{tabularx}{\linewidth}{@{}P{0.38\linewidth}P{0.25\linewidth}Y@{}}
\toprule
\textbf{Synthetic case} & \textbf{Classified verdict} & \textbf{Strong baseline: execution decision}\\
\midrule
W1: local permissions pass; status promotes to order without mandate & Ungoverned effect & Deny\\
P1: valid mandate; realized net outcome $-5$ units & Admissible effect & Permit\\
B1: full enforcement blocks W1 attempt & No effect in $W$ & Deny\\
N1: eligible invocation, no effect despite mandate & No effect in $W$ & Permit\\
U1: effect known, mandate record missing & Non-evaluable guard & Hold\\
F1: effect known, declared evidence insufficient & Ungoverned effect & Deny\\
U2: effect observation missing & Non-evaluable effect & Permit\\
O1: no invocation & Out of scope & Permit\\
\bottomrule
\end{tabularx}
\end{center}

\textbf{Temporal interpretation of the baseline column.} \textsc{Permit}, \textsc{Deny} and \textsc{Hold} are policy-evaluation outcomes \emph{for authorization of the proposed execution at decision time}, not retrospective claims about whether an effect occurred. In U2, an execution may be permitted even though its later effect record is missing. In O1, \textsc{Permit} is a \emph{counterfactual policy verdict on the declared attributes}: no agent invoked the workflow and no live execution-decision event is claimed. W1 has a counterfactual strong-baseline \textsc{Deny} despite its observed synthetic effect because the fully enforcing policy point was absent in that case. The effect classifier and the policy evaluator answer different questions.

\textbf{Reported result (bounded).} All eight declared expected outcomes passed, including the countermodel $L=1,G_e=0$, the control $L=1,G_e=1$, three-valued handling, and the negative-value admissible case. The same W1 defect receives \textsc{Deny} from a complete baseline. This is a \emph{counterexample to the implication from partial local permission to global mandate}, not a counterexample to competent end-to-end authorization, a universal theory, or the published Shadow AI framework. The test is exhaustive only with respect to the eight listed cases, \emph{not} the combinatorial input space or all possible policies.

\textbf{Reproduction contract.} Run \texttt{python3 UT\_v0\_2\_finite\_model.py} (Python $\geq3.9$, no third-party packages). The \emph{unchanged, hash-frozen} v0.2 script and byte-identical reference output \path{UT_v0_2_finite_cases.csv} are supplied in the v0.3.1 reproduction package. The script contains frozen expectations, assertions and the CSV writer. A package README states the execution command, expected stdout and SHA-256 manifest; checksums bind the reproduced artifact to its frozen version. A separate reviewer report communicated an independent successful run under CPython 3.12.14; this reports model reproducibility under the stipulated inputs, not independent validation of its semantics. The classifier has \emph{not} been silently updated to accommodate this manuscript-only notation correction. The encoding does not validate practical guard adequacy or source-to-event attribution.

\section{Credible baseline, reduction and prospective comparison}
The \emph{weak baseline} checks registered agents, service-account access, logging and local endpoint permissions. Under this limited scope W1 passes locally; the order still violates the stipulated cross-boundary mandate. The \emph{strong baseline} checks subject, object, operation, policy version, status meaning, evidence, delegation chain, approver authority and effect-specific scope at the point of execution. It produces \textsc{Deny} for W1 and \textsc{Permit} for P1. Both baselines use the same observation surface, avoiding an unfair information advantage for PADC.

\textbf{Reduction result.} If ordinary ABAC/workflow authorization enforces the complete $G_e$ contract, the proposed PADC guard can be represented as that policy. Within the frozen model, no additional prevention or classification capability has been demonstrated. The distinction between effect and admissibility may still provide explanatory ordering, but any assertion of \emph{superiority} would require independent empirical comparison.

A future diagnostic study should freeze a corpus of cases, anonymized policy artifacts and a gold-standard adjudication procedure \emph{before} allocating reviewers. Compare \textbf{B0} (strong ABAC and workflow/audit review) against \textbf{B1} (identical evidence and controls, with an added PADC transition worksheet). Measure authority-error precision/recall, false positives, inter-rater agreement, adjudication reversals, time and remediation burden. Include matched negative controls and blinded labels. An unchanged or worse B1 performance, at comparable cost and case coverage, defeats the hypothesis of incremental operational usefulness for the tested domain. A study that supplies B1 with additional evidence would \emph{not} isolate vocabulary effects.

\section{Minimal evaluation contract and test protocol}
For empirical use beyond the finite illustration, a case cannot be classified from an evocative narrative alone. Freeze the following interface \emph{before} the assessment:
\begin{enumerate}
\item \textbf{Object and scope:} one AI-mediated transition chain, named systems, principals, action and affected state.
\item \textbf{Window $W$:} starting event and latest time by which an effect would be observable.
\item \textbf{Evidence surface:} output and status records, schema/policy versions, IAM permissions, delegation tokens, workflow rules, approvals, effect logs and independent outcome records.
\item \textbf{Trace rule:} bind each output/status to a subsequent event with stable identifiers; establish whether the effect actually occurred.
\item \textbf{Guard contract:} define stage-specific admission and end-to-end execution mandate, including accepted evidence, permission scope and risk threshold.
\item \textbf{Baseline and counterfactual:} specify whether existing Zero Trust + workflow authorization + audit detects or prevents the same transition without PADC-specific vocabulary.
\item \textbf{Verdict and non-evaluability:} classify only a declared claim; preserve \NE{} for missing or contradictory data.
\end{enumerate}

\begin{center}
\small
\begin{tabularx}{\linewidth}{@{}P{0.32\linewidth}Y@{}}
\toprule
\textbf{Condition under contract} & \textbf{Permitted case verdict}\\
\midrule
Capability shown, no effect observed within adequate $W$ & Non-transduction in $W$; not a permanent impossibility\\
Effect observed, end-to-end guard demonstrably violated & Ungoverned transduction; specify exact broken transition\\
Effect observed, end-to-end guard satisfied & Admissible transduction; value must be evaluated separately\\
Effect or mandate cannot be reconstructed & \NE{}; identify missing evidence and do not infer violation\\
\bottomrule
\end{tabularx}
\end{center}

\Needspace{12\baselineskip}\section{MPRS/ACHO claim inventory and adversarial status}
The internal Michelle Pre-Submission Review System v0.3.2 requires a claim inventory before an applicable adversarial overlay and forbids treating completed formal checks as empirical validation \cite{mprs2026}. This revision distinguishes \emph{executed} model checks from \emph{open} external tests:
\begin{center}
\small
\begin{tabularx}{\linewidth}{@{}P{0.12\linewidth}P{0.42\linewidth}Y@{}}
\toprule
\textbf{ID} & \textbf{Frozen claim and class} & \textbf{Current evidence / status}\\
\midrule
UT-C1 & Window-scoped taxonomy separates no effect, effect with failed guard, and effect with passing guard; definitional. & Explicit decision rule; adequacy remains contestable.\\
UT-C2 & $L=1$ can coexist with $G_e=0$ under incomplete composition contracts; formal possibility. & Constructed witness W1 and executed assertions; model-relative only.\\
UT-C3 & A complete ordinary policy-enforcement baseline detects W1; model-relative comparative finding. & Executed baseline Deny; not a real-deployment claim.\\
UT-C4 & Adding PADC improves diagnosis beyond strong baseline; empirical hypothesis only. & Open; no observed improvement or test corpus.\\
UT-C5 & Admissible effect does not entail positive value; conceptual non-entailment. & Witness P1 with declared $V=-5$; value semantics stipulated.\\
\bottomrule
\end{tabularx}
\end{center}

\textbf{Completed bounded attacks.} Collapse (AH3): UT-C2 reduces to ordinary workflow/attribute authorization when global policy is enforced. Composition (AH5): W1 is the negative witness and P1 is the positive control. Operational instantiation (AH6): the finite script instantiates the classifier and demonstrates the status-promotion bypass only with a deliberately incomplete local implementation. Strong-baseline comparison (AH9): baseline detects W1. Defeat (AH10): specific conditions below. Reproduction (AH11): eight cases executed; this is not exhaustive enumeration. Other ACHO checks remain partial or open: external claim provenance and full-workflow evidence (AH4), predicate separation across real policies (AH7), decision-time observability in deployed systems (AH8), and independent replication. No \emph{No Defeating Finding} certification is claimed.

\textbf{Open adversarial tests / revisions:}
\begin{enumerate}
\item \textbf{Control absorption:} Does complete end-to-end workflow authorization already express every guard that ``ungoverned transduction'' adds? If so, claim conceptual redescription only, not a novel security control.
\item \textbf{Provenance saturation:} If identity, telemetry, provenance, schema mapping, delegated scope and approvals are fully traceable, can a non-admissible promotion still be shown? If not, the strong residual-gap claim fails.
\item \textbf{Scope robustness:} Extend beyond the already executed local/global witness. Which policy bindings or decomposition choices make the counterexample disappear, and which survive legitimate status and role variations?
\item \textbf{Blind-case discipline:} If a reviewer cannot determine the effective authorization chain, does the method return \NE{} instead of inventing an adverse verdict?
\item \textbf{Incremental utility:} On a frozen corpus of synthetic and authorized real cases, can B1 improve authority-error adjudication over B0 given identical evidence? Record false positives, inter-rater agreement, time cost and remediation benefit.
\end{enumerate}

\heading{Explicit defeat condition.} If the proposed PADC analysis provides no reproducible improvement in detecting, classifying or preventing status-to-effect authority errors over an equivalently specified Zero Trust + workflow authorization + audit baseline, the claim of \emph{incremental operational value} must be withdrawn. The conceptual difference between \emph{no effect} and \emph{unauthorized effect} may still be useful, but it is not enough to establish a new mechanism or superior framework.

\section{Limitations, integrity, and research trajectory}
This note offers an analytical taxonomy, one constructed counterexample and eight finite synthetic test cases. It does not estimate the prevalence of ungoverned transitions, establish any causal relation to enterprise profits, validate real-world guard adequacy or show that named frameworks fail when fully implemented. The examples are written to express the thesis and are not an independent sample. Runtime assertions show consistency of the encoding only. ``Value'' is context-sensitive and may include compliance, safety, reliability, time, service quality and costs; any claim about positive value requires its own metric and evaluation window. The boundary between a \emph{missing record} and a \emph{non-compliant record} must remain explicit.

\textbf{Availability and review status.} The complete synthetic-model source, reference CSV, expected console output, reproduction instructions and SHA-256 manifest accompany this manuscript as a separate versioned research package. The package is not an independently maintained public repository, and no DOI is claimed. The note has not undergone journal peer review. Financial-support and competing-interest statements must be checked and approved by the author before public posting or any formal submission.\par

A defensible next research sequence is: (i) external expert review of the definitions and cited sources; (ii) adversarial expansion and independent reproduction of the toy model; (iii) a preregistered, blinded comparison against mature authorization controls; (iv) only then a decision whether the PADC/DEA crosswalk merits publication as a distinct interface artifact. The \emph{Bitter Lesson} essay should remain a separate discussion piece; no causal economic conclusion should be imported from this note.

\section*{Conclusion}
\addcontentsline{toc}{section}{Conclusion}
The model-relative result is that local authorization can coexist with an inadmissible global effect \emph{when} the local contracts leave status-to-mandate promotion unchecked. The complete baseline blocks exactly that case. The proposed distinction is therefore not between ``AI that works'' and ``AI that fails'', but between \emph{effect occurrence} and \emph{warranted transition into effect}. Capability may remain stranded, generate inadmissible effects, or generate admissible effects with positive, neutral or negative value. The researchable question is whether a transition-oriented vocabulary reliably detects authority promotion that otherwise goes unrecognized---and under what conditions that vocabulary is redundant.
\begin{equation*}
\boxed{\text{Effect is not authority; authority is not value.}}
\end{equation*}

\section*{References and provenance}
\addcontentsline{toc}{section}{References and provenance}
\begin{thebibliography}{12}\footnotesize
\bibitem{boehm2026essay} B\"ohm, H. T. (2026). \emph{The Silent AI Bubble: Why Compute Scales but Value Does Not Automatically Transduce}. Discussion Essay v0.2, June 2026. User-supplied LaTeX/PDF; primary conceptual precursor, not peer-reviewed evidence.
\bibitem{shun2026} Shunmugasundaram, M., Gangadharan, S., Dave, R., Kala, S., \& Odikeme Powede Michael (2026). Shadow AI as the new Shadow IT: Governance blind spots in autonomous enterprise systems. \emph{EDPACS}, online 2 June 2026. \url{https://doi.org/10.1080/07366981.2026.2679653}. Used for publisher-verified bibliographic metadata and abstract-level external problem framing; the entire eight-layer design was not independently audited.
\bibitem{rose2020} Rose, S., Borchert, O., Mitchell, S., \& Connelly, S. (2020). \emph{Zero Trust Architecture}. NIST SP 800-207. \url{https://doi.org/10.6028/NIST.SP.800-207}.
\bibitem{hu2019} Hu, V. et al. (2014; updated 2 August 2019). \emph{Guide to Attribute Based Access Control (ABAC) Definition and Considerations}. NIST SP 800-162. \url{https://doi.org/10.6028/NIST.SP.800-162}.
\bibitem{mprs2026} B\"ohm, H. T. (2026). \emph{Michelle Pre-Submission Review System (MPRS) v0.3.2}. Internal candidate revision, 24 October 2026; methodological self-review reference, not external certification.
\bibitem{nist2023} Tabassi, E. (2023). \emph{Artificial Intelligence Risk Management Framework (AI RMF 1.0)}. NIST AI 100-1. \url{https://doi.org/10.6028/NIST.AI.100-1}. Existing governance baseline, not evidence of a PADC-specific deficiency.
\bibitem{padc2026} B\"ohm, H. T. (2026). \emph{Pre-Action Differentiation Cascade (PADC): Source Intake Mode and Epistemic Admissibility}. General Systems Draft v0.4.0, 17 May 2026. User-supplied PDF; project-native guard/intake vocabulary.
\bibitem{interface2026} B\"ohm, H. (2026). \emph{HSRT-Core Interface Specification v0.2.2}. May 2026. User-supplied PDF; application-level object/evidence/trace and non-evaluability discipline.
\bibitem{sutton2019} Sutton, R. (2019). \emph{The Bitter Lesson}. \url{http://www.incompleteideas.net/IncIdeas/BitterLesson.html}. Context for capability scaling, not evidence of effect or governance.
\end{thebibliography}

\vspace{3pt}
\noindent\footnotesize\color{soft}\textbf{Provenance and AI-assistance note (v0.3.1).} Author: Henric T. B\"ohm. Revised using OpenAI ChatGPT for analysis, drafting, executable toy-model construction and LaTeX. Human responsibility and final authorization of publication remain with the author. Source set: supplied \emph{Silent AI Bubble} v0.2, PADC v0.4.0, HSRT Interface v0.2.2, MPRS v0.3.2; externally checked publisher metadata/abstract for Shunmugasundaram et al. and official NIST publication pages (27 October 2026). The full Shadow AI article, independent security review, real system logs, current PADC Control baseline and DEA publication master were not verified. Synthetic cases must not be attributed to those sources. MPRS/ACHO terminology is from the supplied MPRS synthesis of R\'oth-derived review techniques; it is not represented as terminology or endorsement by R\'oth.
\end{document}
